A continuacion se desarrollan los diferentes experimentos llevados a cabo, con la intencion de lograr cumplir los objetivos definidos en la fase de la introduccion \ref{chapt:intro}.

\section{Diseño de la metodología experimental}
% Qué variables quieres estudiar.
% Qué métricas necesitas.
% Qué codificadores vas a comparar.
% Qué características de los vídeos vas a utilizar.
% Qué escenarios vas a estudiar.
% Qué hipótesis planteas.

%PASOS PROCESO EXPERIMENTAL
% Cual es el ciclo de vida de cada experimento
% 
% El inicio del desarrollo del proyecto se ha establecido para el 20 de mayo, coincidiendo con la fecha de inicio de la actividad laboral como Personal No Doctorado Investigador (PDI). De forma simultánea, se iniciará la redacción del Trabajo Fin de Máster (TFM), en la que se irán documentando los avances realizados, así como las decisiones técnicas adoptadas durante cada una de las fases de planificación y desarrollo del proyecto.
% 
% La primera etapa del proyecto consistirá en definir su alcance, los objetivos que se pretenden alcanzar y las hipótesis iniciales, tomando como referencia el estado del arte en el ámbito de la codificación de vídeo y la estimación del consumo energético de cargas de trabajo desplegadas sobre contenedores. Una vez establecidos estos aspectos, se determinarán los requisitos y necesidades técnicas del proyecto y se procederá a la preparación del entorno experimental, incluyendo la instalación y configuración del sistema operativo y de las aplicaciones necesarias para la ejecución de las pruebas.
% 
% A continuación, se realizará la selección de los vídeos que se emplearán en los diferentes experimentos. Para ello, se utilizarán los vídeos disponibles en el servidor de contenido del proyecto Eco-Stream. En concreto, se emplearán los vídeos contenidos en el directorio \textbf{test}, procedentes de CableLabs y Blender \cite{cable-labs} \cite{blender-videos}. Estos vídeos se encuentran divididos en chunks de dos segundos de duración, manteniendo las propiedades de calidad de los vídeos originales.
% 
% Una vez preparado el entorno experimental y seleccionados los vídeos, se llevará a cabo el primero de los experimentos, cuyo objetivo será evaluar la capacidad de las diferentes herramientas de monitorización para estimar el consumo energético de un nodo y de un proceso de codificación ejecutado en un contenedor. Estas pruebas se realizarán en escenarios con múltiples contenedores y diferentes limitaciones de acceso a los núcleos de procesamiento del sistema.
% 
% Tras seleccionar la herramienta de monitorización que se utilizará durante el resto de los experimentos, se procederá a seleccionar un conjunto de chunks de vídeo a partir de los datos proporcionados por el Dr. Francisco Miguel Micó Enguídanos, quien ha empleado estos datos en el desarrollo de su tesis doctoral. Esta selección se complementará con un proceso de selección manual de chunks. Posteriormente, el conjunto seleccionado se clasificará en función de sus valores de Spatial Information (SI) y Temporal Information (TI), permitiendo seleccionar los chunks que se utilizarán para generar los vídeos empleados en las siguientes fases experimentales.
% 
% En el siguiente experimento se analizará el consumo energético asociado a la codificación de vídeo mediante la variación del bitrate y del codificador utilizado, mientras se realizan las correspondientes mediciones de consumo energético. Una vez recopiladas las muestras, se procesarán los resultados y se generarán las gráficas necesarias para su análisis. Finalmente, se redactarán las conclusiones correspondientes a esta fase experimental.
% 
% Por último, se llevará a cabo el último experimento, en el que se estudiará el efecto de variar el número de cores y threads asignados a cada proceso de codificación para los diferentes codificadores. A partir de los resultados obtenidos, se generarán gráficas que permitan analizar la relación entre el consumo energético, el número de cores utilizados y el tiempo requerido para completar cada proceso de codificación. Con base en este análisis, se establecerán las conclusiones correspondientes a esta fase experimental.
% 
% Finalmente, los resultados obtenidos durante la fase experimental serán revisados y contrastados con el tutor del TFM, Juan Gutiérrez Aguado. A partir de esta revisión se establecerán las conclusiones finales del proyecto y se completará la redacción del Trabajo Fin de Máster.
% 
% La estimación temporal de las diferentes actividades descritas se recoge en el diagrama de Gantt presentado a continuación \ref{fig:expect-ganttinicial}, que permite visualizar la duración prevista de cada tarea y su distribución a lo largo del desarrollo del proyecto.

\section{Configuración de Entorno \& Herramientas}\label{exp0-environment}
% - Preparacion del sistema
% - Instalacion de herramientas
% - Resultado final del sistema

% Seleccon de S0 
\subsection{Seleccion e Intalacion del Sistema Operativo}
%DONE
En un principio se decidió instalar el sistema operativo XUbuntu 26.04 en base a la recomendación de mi tutor de TFM, Juan Gutiérrez Aguado, ya que es un sistema operativo que él utiliza diariamente y sobre el que tiene una vasta experiencia de uso. Esto proporciona una capa extra de seguridad y de conocimiento sobre el entorno en el que se desarrollarían las pruebas, además del resto de recursos de conocimiento oficiales y no oficiales que existen sobre este sistema operativo.

Una vez finalizada la instalación, se hizo uso del gestor de versiones de drivers interno de Ubuntu, llamado 'ubuntu-drivers', para instalar la versión de los drivers de NVIDIA más compatible con la tarjeta gráfica NVIDIA Quadro M4000. Al comprobar si se podía acceder a los datos internos de consumo mediante nvidia-smi, se obtuvo un error que indicaba que no era capaz de comunicarse con la API interna de la tarjeta gráfica. Al revisar los mensajes proporcionados por el kernel, se pudo determinar que se trataba de un problema de compatibilidad de los drivers con la versión del kernel. Por este motivo se decidió instalar la versión 24.04 de XUbuntu, la cual hace uso de la versión 6.14 del kernel, solventando el problema de compatibilidad de los drivers.

Una vez reinstalado el sistema operativo, se procedió a reinstalar nuevamente los drivers de NVIDIA junto con los paquetes requeridos por nvidia-smi, de la misma forma que con la versión anterior. De nuevo, al tratar de acceder a los datos de consumo mediante `nvidia-smi`, apareció el mismo problema, pero al revisar los mensajes del kernel, se indicaba que el error no era de compatibilidad, sino la ausencia de un componente concreto llamado 'nvidia-drm'. Este componente forma parte del conjunto de módulos del kernel de NVIDIA que permite al kernel comunicarse con la tarjeta gráfica de NVIDIA \cite{nvidia-drm-documentation}. Al tratar varias veces de reconstruir el módulo manualmente, sin lograr éxito, se optó por buscar información en foros oficiales de NVIDIA y foros de la comunidad Ubuntu y Xubuntu, llegando a la conclusión de que era un problema que sucedía en algunas versiones de XUbuntu debido a dos razones principales: La versión Nouveau de los drivers ,y la incompleta compatibilidad de los drivers con XUbuntu.\

El primer problema estaba relacionado con el controlador libre Nouveau, desarrollado por la comunidad mediante técnicas de ingeniería inversa. Aunque este controlador proporciona compatibilidad básica con un amplio número de tarjetas gráficas NVIDIA, no implementa completamente las funcionalidades propietarias ofrecidas por los controladores oficiales, entre ellas la interfaz NVML, utilizada por herramientas como nvidia-smi para acceder a métricas avanzadas de la GPU, incluido el consumo energético. Para solucionar este problema fue necesario deshabilitar la carga del módulo Nouveau y reinstalar exclusivamente los controladores oficiales proporcionados por NVIDIA.\

El segundo problema estaba relacionado con el mecanismo de instalación automática de controladores proporcionado por la utilidad ubuntu-drivers. Tras consultar la documentación oficial y diferentes foros de las comunidades de Ubuntu y Xubuntu, se comprobó que esta herramienta puede instalar automáticamente la versión más reciente del controlador junto con paquetes pertenecientes a versiones diferentes o dejar dependencias de instalaciones anteriores. Esta situación puede provocar inconsistencias entre los módulos del kernel y las bibliotecas de usuario, impidiendo que componentes como nvidia-drm o nvidia-smi funcionen correctamente. Como consecuencia, fue necesario eliminar completamente todos los paquetes relacionados con NVIDIA y realizar una instalación manual de una versión del controlador compatible con la tarjeta gráfica y con el sistema operativo empleado.\

Pese a todo ello, no se logró una comunicación adecuada con la API interna de la tarjeta gráfica, lo que impidió el acceso a las métricas de consumo energético mediante nvidia-smi. Es por ello que se optó por instalar otra distribución de Linux más flexible, con mayor compatibilidad con los drivers de NVIDIA. Finalmente, se decidió instalar Debian, ya que cumplía con los requerimientos necesarios, además de ser una distribución que a lo largo de los años personalmente había utilizado para diferentes proyectos; por tanto, tenía experiencia de uso y conocía su funcionamiento interno y su filosofía de trabajo.\

%PENDIENTE
La versión de Debian seleccionada una version 2 releases por detras de la version actual (Duke - 15)\cite{debian-releases}, siendo la version 13(Trixie) la seleccionada para la instalacion. Durante la instalacion, se decidio marcar la autodeteccion de tarjeta grafica y la instalacion automatica de drivers de NVIDIA.


% Instalacio de requerimeintos del sistema para pruebas
%   - Nvidia Drivers  
%   - Nvidia SMI
%   - Nvidia Container Toolkit
%   - Docker 
% Seleccion de Herramientas de metricas 
%  - Criterio de Sleccion
%  - Repos Visitados
%  - Papers consultados
% Instalacion y Uso de herramientas
%   - Container de linux-server ffmpeg 
%   - Container de Juan EM
%   - Kepler
%         - Problemas y Modificacion de kepler manual - libnvida-ml
%   - Scaphandre Container

\label{sec:entrono-desarrollo}

\section{Implementación de la infraestructura experimental}
% - Explicar Implementacion de experimentos 
% - Explicacion de los script que he utilizado y su funcionamiento


